%% \exrefcheck with the gb4e syntax, in the preamble, at three words.
%% gb4e sets a plain item straight from the input, so its text is not
%% recorded and the report says so rather than guess; \ex[*]{...} is
%% quoted from its argument, and a gloss from its object tier.  Also the
%% footnote series, by its roman number: refcheck is a PDF/UA case, where
%% a footnote cannot be (see ua.tex), so it is here, in the dot syntax
%% the mixed preamble allows.
\documentclass[11pt]{article}
\usepackage[a4paper,margin=25mm]{geometry}
\usepackage{iftex}
\ifPDFTeX
  \usepackage[T1]{fontenc}
  \usepackage[utf8]{inputenc}
\else
  \usepackage{fontspec}
\fi
\usepackage[lazy,gb4e]{linguexx}
\pagestyle{empty}

\exrefcheck[words=3]
\begin{document}
GBNEXT \Next shows it.

\begin{exe}
\ex GBPLAIN plain item text here.
\ex[*]{GBARG eins zwei drei vier}
\ex \gll GBOBJ Hund bellt laut\\
         GBGL dog barks loudly\\
    \glt `The dog barks loudly.'
\end{exe}

GBLLAST \LLast, GBLAST \Last.

GBFN\footnote{\ex. GBFOOT eins zwei drei vier.

GBFNLAST \Last}
\end{document}
